\documentclass[12pt]{article}

\usepackage[utf8]{inputenc}
\usepackage[T1]{fontenc}

\title{linguistica de la biologia}
\author{Nicol Rubiano, Jireh Gallo}
\date{\today}

\begin{document}

\maketitle

\section{Introducción}
LA BIOLOGIA COMO SISTEMA DE INFORMACION 
En este enayo definiremos la linguistica de la biología, y explicaremos como los organismo vivos no solo son materia y energia,
que deben trasnmitir, traducir e interpretar inormacion

LA SINTAXIS DEL ADN
Abordaremos el codigo genetico desde una perspecyiva linguistica, como la secuencia del ADN funcionan con una sintaxis propia.

COMUNICACION CELULAR
Profundizaremos como las celulas inidividuales actuan como agentes que interpretan su entorno, no solo como coliciones quimica al azar, si mo interprentandolo como un proceso de recepcion y traduccion de significado biologico.

ECOLOGIA Y EVOLUCION
En este apartado llevaremo el analicis linguitico al compo de la biologia y los ecoitemas, los rasgo morologico como la coloracion, de alerta o mimetimo, dejara de entenderse como caracteristica fisica para analizarse como signos con un significado dentro de la red ecologia 



\end{document}
